\documentclass[10pt,a4paper]{amsart}
\usepackage[margin=19mm]{geometry}
\usepackage{amsmath,amssymb,mathtools}
\usepackage[scr=rsfs,cal=euler]{mathalfa}
\usepackage{xfrac}
\usepackage[unicode,hidelinks,bookmarks=false]{hyperref}
\newcommand{\ie}{i.e.\ }
\newcommand{\cf}{cf.\ }
\newcommand{\resp}{respectively\ }
\newcommand{\Subsection}{\subsection}
\providecommand{\nobreakdash}{\nobreak\hbox{-}}
\setlength{\emergencystretch}{2em}
\multlinegap0pt
\usepackage{amssymb}
\usepackage{tikz}
\usepackage{bm}
\usepackage{float}
\newtheorem{proposition}{Proposition}
\newtheorem{theorem}{Theorem}
\title{Generalized Space Time Fractional Skellam Process}
\author{Kartik Tathe, Sayan Ghosh}
\date{Source: arXiv:2504.08374v2 (CC BY 4.0)}
\begin{document}
\maketitle

\noindent\emph{Source and licence.} Generalized Space Time Fractional Skellam Process. Version arXiv:2504.08374v2 (CC BY 4.0), distributed by arXiv under the Creative Commons Attribution 4.0 International licence, http://creativecommons.org/licenses/by/4.0/, as recorded in the arXiv licence metadata for this exact version and shown on the abstract page https://arxiv.org/abs/2504.08374v2. Full licence notice: \texttt{licenses/arxiv-2504.08374-CC-BY-4.0.txt}.\par\smallskip

\noindent\textit{Editorial scope.} Counted unit: the article's theorem recurrence.GSFSP (source0.tex lines 913--918). The article's complete original proof of it is delivered with it (source0.tex lines 920--970, one \texttt{proof} environment of 51 lines). No mathematics is written, repaired or re-derived: the statement and the proof are the article's own lines, in the article's own notation, and no retained line is substituted or re-flowed.  Retained uncounted as the source-local context the proof needs: lines 692--697.  Nothing was omitted from any retained span, so the package carries no omission record.  No retained line contains a citation, so the wrapper carries no bibliography and no bibliography entry.  No retained line contains a figure environment, an image-inclusion command or drawing code, so no image is delivered and the package declares no asset dependency.  Editorial formatting only, in addition: an AMS \texttt{amsart} preamble that restates the article's own declarations and macros used by the retained lines, namely the environments \texttt{proposition}, \texttt{theorem} on separate counters, together with the standard packages those lines need.  The retained environments print in this document as Proposition 1, Theorem 1 in order of appearance, whereas the article prints the same items with the numbering of its own full document.  The complete native source of the article as distributed in the arXiv e-print for this exact version is bundled unchanged as \texttt{sources/176380/source0.tex}.
\begin{proposition}\label{p.g.f._GSTFSP1}
    The p.g.f. of GSTFSP-I for $|u|<1$ is given by
    \begin{equation*}
        \mathcal{G}_{\mathcal{S}^{\alpha_1,\alpha_2}_{\beta_1,\beta_2}}(u,t)=\sum_{i=0}^{\infty}\frac{(-(\sum_{j=1}^{k}(1-u^j)\lambda_j)^{\beta_1} t^{\alpha_1})^i}{\Gamma(i{\alpha_1}+1)}\times\sum_{l=0}^{\infty}\frac{(-(\sum_{j=1}^{k}(1-u^{-j})\mu_j)^{\beta_2} t^{\alpha_2})^l}{\Gamma(l{\alpha_2}+1)}.
    \end{equation*}
\end{proposition}

\begin{theorem}\label{recurrence.GSFSP}
    The state probabilities of GSFSP-I satisfy the following recurrence relation for $|u|<1$, $n\geq 1$:
\begin{equation*}
    p(n,t)=\frac{t}{n}(\beta_1(\sum_{j=1}^{k}(1-u^j)\lambda_j)^{\beta_1-1}\sum_{j=1}^{k}j\lambda_jp(n-j,t)-\beta_2(\sum_{j=1}^{k}(1-u^{-j})\mu_j)^{\beta_2-1}\sum_{j=1}^{k}j\mu_jp(n+j,t)).
\end{equation*}
\end{theorem}

\begin{proof}
Using the definition of p.g.f. for GSFSP-I, we have
\begin{align}
\frac{d}{du} \mathcal{G}_{\mathcal{S}_{\beta_1,\beta_2}}(u,t)&=\frac{d}{du}\sum_{i=0}^{\infty}u^{i}p(i,t)=\sum_{i=0}^{\infty}(i+1)p(i+1,t)u^{i}.\label{dif_p.g.f.1}
\end{align}
From Proposition \ref{p.g.f._GSTFSP1}, for $\alpha_1=\alpha_2=1$, we have 
\begin{align}
\mathcal{G}_{\mathcal{S}_{\beta_1,\beta_2}}(u,t)
    % &=E_{1,1}(-(\sum_{j=1}^{k}(1-u^j)\lambda_j)^{\beta_1} t)\times E_{1,1}(-(\sum_{j=1}^{k}(1-u^{-j})\mu_j)^{\beta_2} t)\nonumber\\
    &=\exp\{-t(\sum_{j=1}^{k}(1-u^j)\lambda_j)^{\beta_1}-t(\sum_{j=1}^{k}(1-u^{-j})\mu_j)^{\beta_2}\}.\label{p.g.f._GSFSP}
\end{align}
Differentiating \eqref{p.g.f._GSFSP}, we get
\begin{align}
    &\frac{d}{du}\mathcal{G}_{\mathcal{S}_{\beta_1,\beta_2}}(u,t)\nonumber\\
    % &=\mathcal{G}_{\mathcal{S}_{\beta_1,\beta_2}}(u,t)(-t)[\beta_1(\sum_{j=1}^{k}(1-u^j)\lambda_j)^{\beta_1-1}(\sum_{j=1}^{k}-ju^{j-1}\lambda_j)+\beta_2(\sum_{j=1}^{k}(1-u^{-j})\mu_j)^{\beta_2-1}(\sum_{j=1}^{k}ju^{-j-1}\mu_j)]\nonumber\\
    &=\sum_{i=0}^{\infty}u^ip(i,t)(-t)[\beta_1(\sum_{j=1}^{k}(1-u^j)\lambda_j)^{\beta_1-1}(\sum_{j=1}^{k}-ju^{j-1}\lambda_j)+\beta_2(\sum_{j=1}^{k}(1-u^{-j})\mu_j)^{\beta_2-1}(\sum_{j=1}^{k}ju^{-j-1}\mu_j)].\label{dif_p.g.f.2}
\end{align}
Comparing the RHS of \eqref{dif_p.g.f.1} and \eqref{dif_p.g.f.2}, we get
\begin{align}
    &\sum_{i=0}^{\infty}(i+1)p(i+1,t)u^{i}\nonumber\\
    &=\sum_{i=0}^{\infty}u^ip(i,t)(-t)[\beta_1(\sum_{j=1}^{k}(1-u^j)\lambda_j)^{\beta_1-1}(\sum_{j=1}^{k}-ju^{j-1}\lambda_j)+\beta_2(\sum_{j=1}^{k}(1-u^{-j})\mu_j)^{\beta_2-1}(\sum_{j=1}^{k}ju^{-j-1}\mu_j)]\nonumber\\
    &=(\sum_{j=1}^{k}(1-u^j)\lambda_j)^{\beta_1-1}\sum_{j=1}^{k}t\beta_1\sum_{i=j-1}^{\infty}j\lambda_jp(i-j+1,t)u^i-(\sum_{j=1}^{k}(1-u^{-j})\mu_j)^{\beta_2-1}\sum_{j=1}^{k}t\beta_2\sum_{i=-j-1}^{\infty}j\mu_jp(i+j+1,t)u^i\nonumber\\
    &=(\sum_{j=1}^{k}(1-u^j)\lambda_j)^{\beta_1-1}\sum_{j=1}^{k}j\lambda_jt\beta_1[\sum_{i=j-1}^{k-1}p(i-j+1,t)u^i+\sum_{i=k}^{\infty}p(i-j+1,t)u^i]\nonumber\\
    &-(\sum_{j=1}^{k}(1-u^{-j})\mu_j)^{\beta_2-1}\sum_{j=1}^{k}j\mu_jt\beta_2[\sum_{i=-j-1}^{k-1}p(i+j+1,t)u^i+\sum_{i=k}^{\infty}p(i+j+1,t)u^i]\nonumber\\
    &=(\sum_{j=1}^{k}(1-u^j)\lambda_j)^{\beta_1-1}\sum_{i=0}^{k-1}\sum_{j=1}^{i+1}j\lambda_jt\beta_1 p(i-j+1,t)u^i+(\sum_{j=1}^{k}(1-u^j)\lambda_j)^{\beta_1-1}\sum_{i=k}^{\infty}\sum_{j=1}^{k}j\lambda_jt\beta_1 p(i-j+1,t)u^i\nonumber\\
    &-(\sum_{j=1}^{k}(1-u^{-j})\mu_j)^{\beta_2-1}\sum_{i=-k-1}^{k-1}\sum_{j=-i-1}^{k}j\mu_jt\beta_2 p(i+j+1,t)u^i-(\sum_{j=1}^{k}(1-u^{-j})\mu_j)^{\beta_2-1}\sum_{i=k}^{\infty}\sum_{j=1}^{k}j\mu_jt\beta_2 p(i+j+1,t)u^i
    \,\,.\label{comparison1}
\end{align}
Now equating the coefficients of $u^{i}$ for $i\leq k-1$ on both sides of \eqref{comparison1}, we obtain
\begin{equation*}
    (i+1)p(i+1,t)=(\sum_{j=1}^{k}(1-u^j)\lambda_j)^{\beta_1-1}\sum_{j=1}^{i+1}j\lambda_jt\beta_1 p(i-j+1,t)-(\sum_{j=1}^{k}(1-u^{-j})\mu_j)^{\beta_2-1}\sum_{j=-i-1}^{k}j\mu_jt\beta_2 p(i+j+1,t)
\end{equation*}
for $-k\leq n\leq k$ which reduces to 
\begin{equation}
    p(n,t)=\frac{t}{n}(\beta_1(\sum_{j=1}^{k}(1-u^j)\lambda_j)^{\beta_1-1}\sum_{j=1}^{i+1}j\lambda_j p(n-j,t)-\beta_2(\sum_{j=1}^{k}(1-u^{-j})\mu_j)^{\beta_2-1}\sum_{j=-i-1}^{k}j\mu_j p(n+j,t))\label{comparison2}
\end{equation}
as in the first summation $p(n-(n+1),t), p(n-(n+2),t),\dots,p(n-k,t)$ are all vanishing and in the second summation $\mu_{-n}, \mu_{-n+1},\dots,\mu_{-1}$ all are zeroes. Also $p(-i,t)=0$ for $1\leq i\leq k$. 

Again, by equating the coefficients of $u^{i}$ for $i\geq k$ on both sides of \eqref{comparison1}, we obtain
\begin{equation*}
    (i+1)p(i+1,t)=(\sum_{j=1}^{k}(1-u^j)\lambda_j)^{\beta_1-1}\sum_{j=1}^{k}j\lambda_jt\beta_1 p(i-j+1,t)-(\sum_{j=1}^{k}(1-u^{-j})\mu_j)^{\beta_2-1}\sum_{j=1}^{k}j\mu_jt\beta_2 p(i+j+1,t)
\end{equation*}
for $n\geq k+1$ which reduces to
\begin{equation}
    p(n,t)=\frac{t}{n}(\beta_1(\sum_{j=1}^{k}(1-u^j)\lambda_j)^{\beta_1-1}\sum_{j=1}^{k}j\lambda_j p(n-j,t)-\beta_2(\sum_{j=1}^{k}(1-u^{-j})\mu_j)^{\beta_2-1}\sum_{j=1}^{k}j\mu_j p(n+j,t)).\label{comparison3}
\end{equation}
Combining \eqref{comparison2} and \eqref{comparison3}, we have for $n\geq 1$
\begin{equation*}
     p(n,t)=\frac{t}{n}(\beta_1(\sum_{j=1}^{k}(1-u^j)\lambda_j)^{\beta_1-1}\sum_{j=1}^{k}j\lambda_jp(n-j,t)-\beta_2(\sum_{j=1}^{k}(1-u^{-j})\mu_j)^{\beta_2-1}\sum_{j=1}^{k}j\mu_jp(n+j,t)).
\end{equation*}
\end{proof}

\end{document}
